% Generated draft: requires mathematical and source review.
\documentclass[11pt]{article}
\usepackage[margin=25mm]{geometry}
\usepackage{fontspec}
\setmainfont{Latin Modern Roman}
\usepackage{amsmath,amssymb,mathtools,unicode-math}
\setmathfont{Latin Modern Math}
\usepackage{enumitem,needspace,xurl,hyperref}
\setlength{\parindent}{0pt}
\setlength{\parskip}{5pt}
\setlength{\emergencystretch}{3em}
\newcommand{\N}{\mathbb{N}}
\newcommand{\Z}{\mathbb{Z}}
\newcommand{\Q}{\mathbb{Q}}
\newcommand{\R}{\mathbb{R}}
\newcommand{\C}{\mathbb{C}}
\newcommand{\F}{\mathbb{F}}
\newcommand{\A}{\mathbb{A}}
\newcommand{\PP}{\mathbb{P}}
\newcommand{\E}{\mathbb{E}}
\newcommand{\eps}{\varepsilon}
\newcommand{\dd}{\,\mathrm{d}}
\begin{document}
\section*{UnsolvedMath ID: 20000233}
\textbf{Source problem code:} catalogue entry
\subsection{Definitions and mathematical statement}
Fix degree $d$ and a generic conjugation-invariant point configuration in $\mathbb P^2$, consisting of $d_1$ real points and $d_2$ nonreal conjugate pairs, with
\[
 d_1+2d_2=3d-1.
\]
The verification field explicitly corrects the extracted ``$2g_2$'' to $2d_2$ and the broken projective-space typography to $\mathbb P^2$. The ambient real rational surface is $\mathbb P^2$; the counted objects are real rational plane curves.

For a nodal enumerated curve $C$, let $s(C)$ be its number of isolated/solitary real nodes, locally having conjugate branches. Its Welschinger weight is $(-1)^{s(C)}$, and the count is
\[
 W_d(d_1,d_2)=\sum_C(-1)^{s(C)}.
\]
The usual genericity, curve class, and real-enumerative conventions must be fixed to make this invariant precise.

\textbf{AIM Problem 9, suggested by Erwan.} Understand these signed counts, particularly their total sign as the mix of real points and conjugate pairs varies. The source remark says formulas exist but their signs are opaque, and suggests degeneration. It reports known signs near either endpoint, nonnegativity when $d_2=0$, and the sign $(-1)^{(d-1)(d-2)/2}$ at minimal $d_1=0$ or $1$, depending on parity. The approximate endpoint ranges are not stated, and the minimal-endpoint sign convention needs clarification when the invariant is zero.

\textbf{Supplied research-aid context, not a verified solution.} Put $\delta_d=(d-1)(d-2)/2$ and let $\Gamma_C$ be the nonisolated real locus of an irreducible real rational nodal degree-$d$ curve, viewed as a nodal graph, not just its smooth normalization. The report claims
\[
 (-1)^{s(C)}=(-1)^{\delta_d-\chi(\Gamma_C)}.
\]
Consequently it rewrites the normalized count as
\[
 (-1)^{\delta_d}W_d=N_{\mathrm{even}}-N_{\mathrm{odd}},
\]
where the two numbers count curves with even or odd Euler characteristic of $\Gamma_C$. It proposes this imbalance as a target for a sign-reversing injection but supplies no such general injection.

For degree three, the attempt considers a generic pencil through eight imposed conjugation-invariant points, of which $2a$ are real. If it has $2t$ real singular members, the report claims the hyperbolic and elliptic nodal counts are
\[
 (H,E)=(t+a,t-a),\qquad
 0\le a\le4,\quad a\le t\le6,
\]
and every such integer profile is realizable. These are the imposed eight points, not a claim that a cubic pencil has only eight geometric base points. Hyperbolic nodes give weight $+1$ and elliptic solitary nodes weight $-1$, so this reported profile gives $W_3=H-E=2a$, including zero at $a=0$. Thus a literal nonzero ``sign'' reading of the original minimal-endpoint remark must not be imposed without a zero convention. The full realizability theorem and parity calculation are unreviewed partial claims, not a sign theorem for all degrees.
\subsection{Short English statement}
Study the sign of plane Welschinger counts with mixed real and conjugate-point constraints. The attached Euler-parity reformulation and cubic cancellation profiles are partial unreviewed claims, with zero-count and endpoint conventions kept explicit.
\subsection{Sources}
\begin{itemize}
\item Local UnsolvedMath registry, permanent ID 20000233, source code catalogue entry; the source says ``Suggested by Erwan,'' without a surname or proposal year.
\item AIM workshop \emph{Degenerations in Algebraic Geometry}, Problem 9 and Remark 9.1; supplied locator \url{https://aimath.org/pastworkshops/degenalggeomproblems.pdf}. Canonical locator: aim-algebraic-geometry-notes.json, notes[232], tag ``problem.'' The verification field identifies the two material extraction errors. Original literature field is empty. No PDF was fetched or generated here.
\item Research attempt 1: partial\_result, status at run ``partially\_solved,'' machine-assigned rigor ``proved,'' novelty confidence ``low,'' with an explicit machine-generated/not-peer-reviewed/verify-before-citing caveat. No exact primary citation for the formulas or cubic realizability theorem is supplied.
\end{itemize}
\end{document}
